High performance computing (HPC) data management systems, i.e., parallel file systems, are treated as I/O devices rather than data movement orchestrators for application workflows, leaving HPC users to manage this orchestration. As these workflows grow in complexity and scale, the burden of data movement increases on users. 
To relieve users from this responsibility, we propose Data Flyway, a data management runtime system for performing ``in-flight processing'', where transformations {\color{red}and analysis operations} are performed on data while it is transferred from source to destination. We present a novel approach for representing data transformation {\color{red}and analysis} pipelines that can be registered with I/O resources, enabling data to be processed transparently and asynchronously during data transfers. We use directed graphs as representations of the in-flight processing pipelines, allowing the system to select the optimal path at runtime based on current resource availability and data movement costs. The directed graph can then be associated with various I/O resources, such that any subsequent read or write operation on that resource automatically triggers the appropriate transformations {\color{red}or analysis operations}. Our strategy decouples transformation {\color{red}and analysis} management from application logic, freeing developers from having to explicitly invoke transformations {\color{red}or analysis operations} on every read or write. Subsequent reads and writes flow through the graph, transforming the data to the desired state {\color{red}or deriving new data} before reaching its destination. We implement this in-flight processing approach in a runtime data management library called Proactive Data Containers (PDC) and demonstrate that this strategy can significantly reduce the complexity of managing data-intensive workloads while ensuring efficient I/O.